\documentclass[a4paper,10pt]{article}

\usepackage[utf8]{inputenc}
\usepackage{tabularx}
\usepackage{csquotes}
% \usepackage[none]{hyphenat}
\usepackage[french=nohyphenation]{hyphsubst}
\usepackage{graphicx}
\usepackage{hyperref}
\usepackage{fancyhdr}
\usepackage{float}
\usepackage{caption}
\usepackage{listings}
\usepackage{spverbatim}
\usepackage{xcolor}
\usepackage[backend=bibtex]{biblatex}
\usepackage[toc]{glossaries}
\usepackage{wrapfig}
\usepackage[french]{babel}
\usepackage{booktabs}

\captionsetup{justification=centering}

\newcommand{\heading}[1]{\multicolumn{1}{c}{#1}}

\lstdefinestyle{customc}{
  belowcaptionskip=1\baselineskip,
  breaklines=true,
  frame=L,
  xleftmargin=\parindent,
  language=C,
  showstringspaces=false,
  basicstyle=\footnotesize\ttfamily,
  keywordstyle=\bfseries\color{green!40!black},
  commentstyle=\itshape\color{purple!40!black},
  identifierstyle=\color{blue},
  stringstyle=\color{orange},
}

\lstdefinestyle{customasm}{
  belowcaptionskip=1\baselineskip,
  frame=L,
  xleftmargin=\parindent,
  language=[x86masm]Assembler,
  basicstyle=\footnotesize\ttfamily,
  commentstyle=\itshape\color{purple!40!black},
}

\lstdefinestyle{custombash}{
  belowcaptionskip=1\baselineskip,
  breaklines=true,
  frame=L,
  xleftmargin=\parindent,
  language=bash,
  showstringspaces=false,
  basicstyle=\footnotesize\ttfamily,
}

\lstset{escapechar=@}

\rhead{BAUDVIN CASSAGNES PUISSANT}
\chead{}
\lhead{NetFLow et ALCASAR}

\pagestyle{fancy}

\sloppy

\addbibresource{biblio.bib}
% \makeglossary
% % \newglossaryentry{<label>}{<settings>}

\newglossaryentry{NAC}{
  name=NAC,
  description={Network Acces Control en anglais est un contrôleur d'accès. Il va permettre à une entreprise de soumettre l'accès à un réseau informatique à un système d'authentification de l'utilisateur}
}
\newacronym{alca}{ALCASAR}{Application Libre pour le Contrôle d'Accès Sécurisé et Authentifié au Réseau}



\begin{document}

%\maketitle
\begin{titlepage}
      \begin{center}   
        \Huge
        \textbf{NetFlow et ALCASAR}
        
        %\vspace{0.5cm}
        \LARGE
        ~
        
%         \vspace{0.5cm}
        
        \vfill
        \begin{figure}[H]
	    \centering
	    \begin{minipage}{0.49\textwidth}
		\centering
		\includegraphics[width=\textwidth]{./img/esiea.jpeg}
	    \end{minipage}
	    \begin{minipage}{0.49\textwidth}
		\centering
		\includegraphics[width=.8\textwidth]{./img/Alcasar-logo.png}
	    \end{minipage}\hfill
	\end{figure}
        \vfill
        
        \vspace{0.5cm}
        
        Analyse de flux IP avec le protocole NetFLow
        
        \vspace{2cm}
        \textbf{Emmanuel BAUDVIN\\Lucien CASSAGNES\\Antoine PUISSANT}\\
        \vspace{0.8cm}
        \Large
        \underline{Enseignant :} M. REY\\
        \vspace{0.5cm}
        2015 - 2016%\today
        
    \end{center}
\end{titlepage}

\begin{abstract}
ALCASAR\footnote{Application Libre pour le Contrôle d'Accès Sécurisé et Authentifié au Réseau} est un projet libre de contrôle d'accès au réseau, un NAC\footnote{Network Access Control}. Au sein de ce dernier, il est possible de retrouver l'utilisation du protocole NetFlow. ALCASAR bénéficie alors d'une sonde réseau lui permettant d'analyser les flux réseaux le traversant avec beaucoup de finesse.\\
Durant ce projet, le but est d'assimiler le fonctionnement du protocole NetFlow seul et au sein d'ALCASAR.\\
Une fois cet apprentissage fait, nous devrons comparer la version de NetFlow présente dans ALCASAR à la dernière version. Cette comparaison devra permettre de savoir s'il est utile de faire évoluer la version d'ALCASAR. En cas de résultats positifs, la mise à jour devra être réalisée et testée.
\end{abstract}

\newpage

\tableofcontents

\section{Introduction}
\subsection{Le projet ALCASAR}

\begin{wrapfigure}{r}{0.30\textwidth}
  \vspace{-20pt}
  \begin{center}
    \includegraphics[width=0.25\textwidth]{img/Alcasar-logo.png}
  \end{center}
  \vspace{-15pt}
  \caption{Logo d'ALCASAR}
  \vspace{-10pt}
\end{wrapfigure}
Le projet ALCASAR est né suite à une demande d'une entité gouvernementale voulant assurer la traçabilité et l'imputabilité des connexions de tous les utilisateurs connectés à leur réseau.\\
Après une rapide étude de marché, les responsables de projet se sont rendu compte qu'il n'existait pas de solution libre de droits permettant de remplir les demandes précises de l'entité : intercepter, authentifier, filtrer et imputer l'accès aux utilisateurs à internet.\\
Ainsi, les chefs du projet ont alors décidé de réaliser eux-mêmes un contrôleur d'accès pouvant répondre aux besoins du projet. C'est ainsi que le projet ALCASAR à vu le jour.\\\par
Afin de pouvoir réaliser ces fonctions d'interception ou de filtrage, ALCASAR va embarquer plusieurs outils d'analyse de flux. Parmi ces outils, nous allons pouvoir retrouver NetFlow.

\subsection{NetFlow}

NetFlow est un protocole qui a été conçu par Cisco Systems. Il va permettre de réaliser une surveillance des réseaux en collectant des informations sur les flux IP. Ce produit va permettre à un administrateur réseau de déterminer différentes informations telles que la source et la destination du trafic, le type de service utilisé ou encore la cause de nœuds de congestion.

\subsection{IPFIX}
A voir : \url{http://www.bradreese.com/blog/netflow-vs-ipfix-exporter.htm}
\section{NetFlow}
\subsection{Indtoduction}
NetFlow est un protocole développé par la société Cisco dans le but de faire une analyse des trames passant au travers de leurs équipements réseau. Ce projet ayant une très grande popularité lors de sa sortie, le protocole devint alors accessible à tous.\\
Ainsi, ici, afin d'utiliser le protocole NetFlow, nous pouvons retrouver plusieurs parties :
\subsubsection{La sonde}
La mise en place du protocole NetFlow est ici faite grâce un module noyau, \textit{ipt\_NETFLOW}.\\
TODO\\
Cette sonde va envoyer toutes ses données à une adresse IP et un port bien spécifique (ici, 127.0.0.1:2055).
\subsubsection{Le collecteur}
Afin de récolter toutes les données fournies par la sonde NetFlow, nous devons aussi avoir un collecteur. Son rôle est d'écouter à l'adresse sur laquelle la sonde NetFlow envoie ses données. Il va ainsi réaliser un fichier au format \enquote{\textit{NetFlow}} toutes les 5 minutes.\\
Ici, le rôle du collecteur est réalisé par le daemon \textit{Nfcapd}.
\subsubsection{L'interpréteur}
Cependant, les données fournies par Nfcapd au format \enquote{\textit{NetFlow}} ne sont pas lisible. Afin de palier à cela, il est possible d'utiliser un interpréteur. Dans notre cas, ce rôle est joué par \textit{Nfdump}.\\
Il est alors possible d'avoir le contenu des captures de la sonde NetFlow à n'importe quel moment.
\subsubsection{Le grapheur}
Cependant, au sein de l'interface de gestion d'ALCASAR (ACC\footnote{ALCASAR Control Center}), il est intéressant de pouvoir retrouver de manière graphique les résultats des captures de la sonde NetFlow.\\
Afin de réaliser cela, nous utilisons le grapheur Nfsen. Ce dernier va récupérer les données fournies par le collecteur afin de réaliser des graphes représentatifs de l'état des connexion (charges, nombre de connexion, serveurs les plus demandés, etc\ldots).

\subsection{Mise à jour}
Afin de mettre toute la partie NetFlow à jour au sein d'ALCASAR, nous avons tout d'abord chercher quelles étaient les versions installées sur la dernière version du NAC.
\begin{table}[H]
    \centering
    \begin{tabularx}{\textwidth}{|l|c|c|}
	\hline
	& \textbf{Version au sein d'ALCASAR} & \textbf{Dernière version}\\\toprule
	\textbf{ipt\_NETFLOW} & 1.7.2 & 2.2-36\\\hline
 	\textbf{Nfdump} & 1.6.9 & 1.6.14\\\hline
 	\textbf{Nfsen} & 1.3.7 & 1.3.7\\\hline
 	\textbf{Nfcapd} & 1.6.9 & 1.6.14\\\hline
    \end{tabularx}
    \caption{Comparaison des versions des modules}
    \label{tab:1}
\end{table}
Comme vous pouvons le voir sur le tableau \ref{tab:1}, nous devons mettre à jour la sonde ipt\_NETFLOW, Nfdump ainsi ue Nfcapd.\\
Actuellement, Nfcapd est inclut au sein de Nfdump. Ainsi, nous devons mettre à jour la sonde et l'interpréteur.\\
Pour cela, nous avons récupéré les sources des deux paquets sur leurs sources officielles :
\begin{description}
 \item[ipt\_NETFLOW] \url{https://github.com/aabc/ipt-netflow}
 \item[Nfdump] \url{https://github.com/phaag/nfdump}
\end{description}

\subsubsection{ipt\_NETFLOW}
Afin d'installer la sonde sur notre poste ALCASAR, nous devons ajouté ce module au noyau. Cependant, le noyau d'ALCASAR est déjà configuré pour la mise en place d'un module NetFlow.\\
ipt\_NETFLOW étant un module supplémentaire d'iptables, il faut s'assurer que ce dernier comporte bien toutes les librairies nécessaires. Pour ce faire, il est recommander de télécharger et installer iptables depuis les sources\footnote{\url{ftp://ftp.netfilter.org/pub/iptables/}}. Il faut cependant faire attention à bien prendre la même version que celle installé de base dans le noyau de la machine. Dans notre cas, pour Mageia 5, la version d'iptables en question est la 1.4.21\footnote{\url{ftp://ftp.netfilter.org/pub/iptables/iptables-1.4.21.tar.bz2}}.\\
Une fois le paquet téléchargé, nous pouvons décompresser l'archive et l'installer :
\begin{lstlisting}[style=custombash, language=bash]
cd iptables-1.4.21
./configure --prefix=/usr/
make
make install 
\end{lstlisting}
Dans certains cas, il est possible de vouloir utiliser un système de supervision ou d'hypervision dans le parc informatique où est installé ALCASAR. C'est dans cette optique que les développeurs d'ipt\_NETFLOW ont pensé à intégrer le module net-snmp. \textbf{\textit{DESCRIPTION DU MODULE}}.\\
L'installation de ce module n'est pas obligatoire pour le bon fonctionnement de NetFlow. Nous avons cependant préféré tester son installation afin de fournir une installation complète de NetFlow.\\
Il est possible de télécharger la dernière version de ce module sur le site officiel de net-snmp\footnote{\url{www.net-snmp.org}}. Dans notre cas, il s'agit de la version 5.3.7.\\
Il faut aussi noter que pour que l'installation de ce module fonctionne, il est nécessaire d'avoir des \textit{header} pour perl. Ces fichiers n'étant pas présents lors de nos tests sur Mageia 5, nous avons du installer le paquet \textit{perl-devel}.
\newpage
Une fois téléchargé, nous pouvons l'installer :
\begin{lstlisting}[language=bash,style=custombash]
cd net-snmp-5.7.3
./configure --prefix=/usr/
make
make install
\end{lstlisting}
Une fois ce module installé correctement, il est enfin possible de réaliser l'installation d'ipt\_NETFLOW. La version utilisée est la dernière version disponible du le GitHub du projet (2.1.36-g95fabbf). Git n'étant pas installé de base sur Mageia 5, nous fournissons l'archive comprenant cette version\footnote{ipt-netflow22.tar.gz}.\\
Si vous choisissez de cloner le dépôt vous même, prenez soin de regarder que la version est la bonne: lancez le script \textit{version.sh} (dans l'archive fournie la version de la branche master est 2.1.36-g95fabbf et non 2.1 !)\\
Une fois décompressée, il est possible de l'installer de la manière suivante :
\begin{lstlisting}[language=bash,style=custombash]
 cd ipt-netflow22
 ./configure
 make all install
 depmod
\end{lstlisting}
Une fois l'installation réalisée, nous pouvons redémarrer la machine afin que le script de démarrage d'ALCASAR initialise correctement la sonde NetFlow.

\subsubsection{Nfdump et Nfcapd}
Pour mettre à jour Nfdump et Nfcapd, il faut d'abord installer les dépendances suivantes:
\begin{itemize}
\item Flex (ou Lex): dans sa dernière version (pris en charge par le script de mise à jour)
\item Bison (ou Yacc): dans sa dernière version (pris en charge par le script de mise à jour)
\item M4: dans sa dernière version (pris en charge pas le script de mise à jour). C'est un language de traitement de macro pour l'installation de nfdump.
\item 
\end{itemize}
Pour chacune des dépendances (si le script n'est pas utilisé) il faut effectuer les commandes d'installation basiques:
\begin{lstlisting}[language=bash,style=custombash]
./configure 
make
make install
\end{lstlisting}
Le ./configure peut être aggémanté de l'option --prefix=/usr/ car de base, les librairies vont être installées dans le dossier lib/ de /usr/local/ au lieu de /usr/lib/ comme cela a été fait pour Alcasar.\\
Une fois les dépendances installées, il suffit d'aller chercher la dernière version de nfdump et de l'installer. Comme auparavant, une fois l'archive décompressée, il suffit de faire les commandes suivantes
\begin{lstlisting}[language=bash,style=custombash]
./configure --prefix=/usr/
make
make install
\end{lstlisting}
Pour vérifier que la version de nfdump a bien été mise à jour, veuillez vérifier que la commande nfdump -V renvoie bien la version que vous voullez installer.


\subsection{Le Script}
Dans le cas où le script d'installation serait perdu, nous fournissons le script qui a été utilisé ici.\\
Si l'archive de ipt-tables est aussi perdue, il suffit de faire un
\begin{lstlisting}[language=bash,style=custombash]
 git clone git://github.com/aabc/ipt-netflow.git ipt-netflow22
\end{lstlisting} 
Cependant si des modifications sont faites dans le code entre le temps où se rapport est rédigé et le temps où vous le lisez, il est fort probables qu'il y ai eu des modifications dans les sources qui peuvent potentiellement ajouter des erreurs dans l'installation de ipt-netflow (si aucune version de release a été ajoutée).

\begin{lstlisting}[language=bash,style=custombash]
#!/bin/sh

INSTALL_REP=TMP_REP

mkdir $INSTALL_REP
cd $INSTALL_REP
wget http://downloads.sourceforge.net/project/nfdump/stable/nfdump-1.6.13/nfdump-1.6.13.tar.gz
wget http://downloads.sourceforge.net/project/flex/flex-2.6.0.tar.gz
wget ftp://ftp.gnu.org/gnu/m4/m4-1.4.17.tar.gz
wget ftp://ftp.gnu.org/gnu/bison/bison-3.0.4.tar.gz
wget http://sourceforge.net/projects/net-snmp/files/net-snmp/5.7.3/net-snmp-5.7.3.tar.gz
wget ftp://ftp.netfilter.org/pub/iptables/iptables-1.4.21.tar.bz2

cp ipt-netflow22.tar.gz $INSTALL_REP

for f in *.tar.*
do 
	tar -xvf $f
done

echo "Updating collector (nfdump) to version 1.6.13"

#INSTALL M4
cd m4-1.4.17
./configure
make
make install
cd ..

#INSTALL FLEX
cd flex-2.6.0
./configure
make
make install
cd ..

#INSTALL BISON
cd bison-3.0.4
./configure
make
make install
cd ..

#INSTALL NFDUMP
cd nfdump-1.6.13
./configure --prefix=/usr/
make prefix=/usr/
make install prefix=/usr/
cd ..

echo "Updating ipt-netflow to version 2.1.36-g95fabbf"

#INSTALL IPTABLES
cd iptables-1.4.21
./configure --prefix=/usr/
make
make install
cd ..

#INSTALL SNMP
urpmi perl-devel
cd net-snmp-5.7.3
./configure --prefix=/usr/
make
make install
cd ..

#INSTALL netflow2.2
cd ipt-netflow22
./configure
make all install
depmod
cd ..

cd ..
rm -rf $INSTALL_REP

nfdump -V

\end{lstlisting} 

\newpage
\rhead{2014-2015}
\nocite{*}
\printbibliography\addcontentsline{toc}{section}{Références}

%\listoffigures\addcontentsline{toc}{section}{\listfigurename}
\end{document}
